% Member 4 report sections: platform (HLD 4 architecture, HLD 5 services 1 and 10, HLD 8 API and events,
% HLD 10 security, HLD 11 tech stack, HLD 12 monitoring, one HLD 13 trade-off).
% Body only, so the team report can % Member 4 report sections: platform (HLD 4 architecture, HLD 5 services 1 and 10, HLD 8 API and events,
% HLD 10 security, HLD 11 tech stack, HLD 12 monitoring, one HLD 13 trade-off).
% Body only, so the team report can % Member 4 report sections: platform (HLD 4 architecture, HLD 5 services 1 and 10, HLD 8 API and events,
% HLD 10 security, HLD 11 tech stack, HLD 12 monitoring, one HLD 13 trade-off).
% Body only, so the team report can % Member 4 report sections: platform (HLD 4 architecture, HLD 5 services 1 and 10, HLD 8 API and events,
% HLD 10 security, HLD 11 tech stack, HLD 12 monitoring, one HLD 13 trade-off).
% Body only, so the team report can \input{member4-platform}. Needs: array, booktabs.
% member4-standalone.tex compiles these sections on their own. Numbers measured on main (b39692e), 21 November 2026.

\section{Architecture (HLD 4)}

The prototype is a \emph{modular monolith}: one FastAPI process that runs all eleven HLD services as eleven Python
packages under \texttt{app/modules/}. HLD 13.1 recommends this for a student build, and it keeps the HLD's service
boundaries, so each package can later move to its own process without a redesign. Every package has the same two
files: \texttt{service.py} holds the business logic as plain functions whose first argument is the database session,
and \texttt{router.py} holds the HTTP endpoints, which parse the input, call one service function and shape the
answer. Because the logic never touches HTTP, modules call each other directly and tests can call the logic without
a server.

Three rules hold the system together:
\begin{itemize}
  \item \textbf{One owner for stock.} Only the inventory module changes stock quantities: what is on hand, what is
    held and which bin it is in. Other modules ask it to hold, release, sell, receive or write off stock; forecasting
    and replenishment only set an item's reorder settings.
  \item \textbf{Hold at checkout, sell at packing.} Checkout raises \texttt{qty\_reserved}; nothing leaves
    \texttt{qty\_on\_hand} until the order is packed. A failed payment or a cancel only gives a hold back, so no stock
    count ever has to be ``put back''.
  \item \textbf{Check and change in one statement.} Every contested change is one conditional \texttt{UPDATE} with
    the check in its \texttt{WHERE} clause, so two requests can never both pass a check that only one of them should.
    The checkout hold, for example:
    \begin{center}\small
      \verb|UPDATE stock_levels SET qty_reserved = qty_reserved + :qty|\\
      \verb|WHERE store_id = :store AND sku_id = :sku AND qty_on_hand - qty_reserved >= :qty|
    \end{center}
\end{itemize}

Table~\ref{tab:hld-vs-prototype} lists where the prototype stands in for an HLD component.

\begin{table}[ht]
  \centering
  \footnotesize
  \begin{tabular}{@{}>{\raggedright\arraybackslash}p{0.2\linewidth}>{\raggedright\arraybackslash}p{0.44\linewidth}
                  >{\raggedright\arraybackslash}p{0.3\linewidth}@{}}
    \toprule
    HLD component & In the prototype & What a full build adds \\
    \midrule
    API gateway & The FastAPI app itself: login, roles, own-store checks, audit log, rate limit & A separate gateway
      with TLS and mutual TLS between services \\
    \addlinespace[3pt]
    PostgreSQL, sharded by store & One SQLite file (WAL mode), all stores in one database & PostgreSQL and
      migrations (task 1.7); sharding by \texttt{store\_id} \\
    \addlinespace[3pt]
    Redis ATP cache & A 2-second in-process cache of each store's availability list, cleared on every stock write
      & A shared Redis, so several app instances see one copy \\
    \addlinespace[3pt]
    Kafka & A transactional outbox (\texttt{event\_log}) and a relay inside the app that feeds each store's live feed
      & A Kafka producer as one more relay sink (section~\ref{sec:m4-tradeoff}) \\
    \addlinespace[3pt]
    Store edge node & A design note and an outage experiment (\texttt{docs/EDGE\_NODE.md}) & A second copy of the
      app in each store, syncing through the same relay \\
    \addlinespace[3pt]
    Prometheus, Grafana & Both, in Docker Compose, with the HLD 12 alert rules & An Alertmanager to send the
      alerts somewhere \\
    \bottomrule
  \end{tabular}
  \caption{HLD components and their stand-ins in the prototype.}
  \label{tab:hld-vs-prototype}
\end{table}

\section{Serviceability (HLD 5, service 1)}

Serviceability answers ``which store delivers here, and how fast?'' for every address a customer types, so it must
be cheap. Delivery zones are H3 hexagon cells at resolution 9: cells with edges of about 0.2\,km, so a zone border is
accurate to about 200\,m and a 3\,km zone is about 270 cells.
\begin{itemize}
  \item \textbf{Zones are built once.} \texttt{rebuild\_zones} gives each store every cell whose centre lies within its
    radius; where two stores claim a cell, the nearer one wins. The result is one row per cell.
  \item \textbf{A lookup is one row.} \texttt{check} turns the customer's point into its cell and reads that cell's
    row, instead of measuring the distance to every store.
  \item \textbf{Load shapes the answer.} A closed store serves nobody. A busy store (more than 5 unpacked orders)
    serves only the inner 70\% of its radius and quotes 5 minutes more. The ETA is 2 minutes to pack plus 3 minutes
    per kilometre, rounded up.
  \item \textbf{A miss is recorded.} A point no store serves is saved as its cell and the time, nothing else. These
    misses are the demand that network planning works from.
\end{itemize}
Under load (section~\ref{sec:m4-monitoring}) a serviceability check took 10.7\,ms at the median and 26.1\,ms at p99.

\section{Network Planning (HLD 5, service 10)}

Network Planning answers ``where should the next dark store go?'' (FR-12). It counts demand per H3 cell, never as a
name, phone number or exact address: a paid order carries the cell of its delivery address, and a turned-away
lookup is a miss. \texttt{GET /v1/network-planning/heatmap} returns both counts per cell over the last 30 days, and
the map at \texttt{/network.html} draws them over the current zones.

\texttt{suggest\_sites} chooses up to $k$ new sites by \emph{greedy maximum coverage} of the unmet demand. The
candidates are the cells with misses, except cells inside an existing zone and the cells listed in
\texttt{NO\_GO\_CELLS} (standing in for the zoning check of HLD 13.2). Each round picks the candidate whose 3\,km
radius covers the most unmet demand, breaking ties towards the more central site, then removes that demand and
repeats. It stops early when no candidate would gain anything. Greedy maximum coverage is guaranteed to cover at
least $1 - 1/e \approx 63\%$ of what the best choice of $k$ sites from the same candidates would cover. Only the 300
busiest miss cells are tried as sites, which bounds the work however many misses pile up.

\section{API and events (HLD 8)}

The API has 63 operations on 61 paths, all versioned under \texttt{/v1/} except \texttt{/health}, and documented by
OpenAPI at \texttt{/docs}.
Table~\ref{tab:api} groups them by service. They follow four conventions:
\begin{itemize}
  \item \textbf{One error shape.} Every refusal is \verb|{"error": {"code", "message", "details"}}|, with 409 when
    the request is valid but the current state forbids it (not enough stock, wrong job state), 422 for bad input and
    404 for something that does not exist. The code is stable, so a screen can react to it.
  \item \textbf{Order steps are safe to repeat.} Paying, packing, cancelling and handing over check the current state
    in the same conditional update, so a repeated call reports the state instead of acting twice. Stock adjustments
    take an \texttt{Idempotency-Key} header. A repeated hold or scan is not deduplicated: it holds or scans again.
  \item \textbf{Staff calls carry a token}; the customer calls (serviceability, availability, checkout) are open but
    rate-limited (section~\ref{sec:m4-security}).
  \item \textbf{Every response has a request id} (\texttt{X-Request-ID}), which also appears on every log line made
    while the request was handled.
\end{itemize}

\begin{table}[ht]
  \centering
  \small
  \begin{tabular}{@{}lrl@{}}
    \toprule
    Service & Operations & Examples \\
    \midrule
    Inventory           & 19 & availability, holds, adjustments, cycle counts, near-expiry \\
    Fulfilment          &  8 & cancel, handover, rider check-in, metrics, attention list \\
    Pick \& Pack        &  5 & next order or batch, scan, empty shelf, pack \\
    Replenishment       &  5 & run, dispatch, reorder settings \\
    Assortment          &  5 & catalogue, recommendations \\
    Forecasting         &  4 & run, compare models \\
    Dashboard           &  4 & store dashboard, city overview, KPIs, events \\
    Serviceability      &  3 & check, stores, rebuild zones \\
    Network planning    &  3 & heatmap, suggestions, coverage \\
    Admin               &  3 & reset demo data, audit log, health check \\
    Slotting, inbound, login & 4 & bins, suggest a bin, goods receipt, login \\
    \bottomrule
  \end{tabular}
  \caption{API operations by service.}
  \label{tab:api}
\end{table}

\paragraph{Events.} The code publishes 17 event types, such as \texttt{order.paid} when a customer pays and
\texttt{rider.assigned} when a bag is packed. Publishing writes the event to the \texttt{event\_log} table
\emph{inside the caller's transaction} and calls the subscribed handlers in that same transaction: this is
the transactional outbox of HLD 9.2. If the stock change rolls back, the event never existed;
if it commits, the event is stored. A relay inside the app reads unsent events every 0.2\,s, oldest first, up to 100
at a time, and hands each one to every sink. Today the only sink is the live feed: a WebSocket per store, so every
screen of that store sees an order paid or a line dropped within a second, without polling. An event is marked sent
only after every sink took it, so delivery is \emph{at least once}, and a consumer must ignore an event id it has
already seen.

\section{Security (HLD 10)}
\label{sec:m4-security}

\begin{itemize}
  \item \textbf{Passwords} are stored as salted scrypt hashes and compared in constant time; none is stored as typed.
  \item \textbf{Login} returns a signed token (JWT, HS256) that lasts 30 minutes and names the user, their role and
    their store. Six roles exist: picker, packer, receiver, manager, ops and admin.
  \item \textbf{Every staff endpoint} checks the role and, for store data, that the store is the caller's own:
    a picker of store 1 gets 403 for store 2. Ops may read every store, and \texttt{/v1/admin/} needs admin.
  \item \textbf{An audit log} keeps every call a logged-in user makes to a protected endpoint, refused or not: who,
    what (method, route and path), the status they got and the request id. An admin reads it at
    \texttt{GET /v1/audit}.
  \item \textbf{The open customer endpoints are rate-limited} per client address by a token bucket: 20 calls a second
    with bursts of 40, then 429 with a \texttt{Retry-After} header. This also limits how fast a caller can flood
    network planning with made-up misses.
  \item \textbf{Privacy.} Demand is kept as H3 cells, about 0.1\,km$^2$ each, so the planning data names no customer
    and no exact address. \texttt{/metrics} can be closed with \texttt{METRICS\_TOKEN}.
\end{itemize}
The limits are listed in Table~\ref{tab:m4-limits}: the built-in signing key and the shared demo password are public
for local development (the app warns about the key at start), a token cannot be revoked before it expires, and the
rate limit is
counted per process.

\section{Tech stack (HLD 11)}

\begin{table}[ht]
  \centering
  \footnotesize
  \begin{tabular}{@{}>{\raggedright\arraybackslash}p{0.22\linewidth}>{\raggedright\arraybackslash}p{0.4\linewidth}
                  >{\raggedright\arraybackslash}p{0.32\linewidth}@{}}
    \toprule
    Layer & Used & Why \\
    \midrule
    Language, web & Python 3.12, FastAPI, Uvicorn, Pydantic & Async server, validation and OpenAPI docs for free \\
    Data & SQLAlchemy 2 on SQLite (WAL) & Runs anywhere in one file; the same ORM code moves to PostgreSQL \\
    Geography & H3 (\texttt{h3} 4) & Equal-area hexagons make zones and demand one indexed lookup \\
    Login & PyJWT, scrypt from the standard library & No extra service for a prototype \\
    Monitoring & \texttt{prometheus-client}, Prometheus 2.53, Grafana 11.1 & The HLD's choice; one Compose profile \\
    Screens & HTML, CSS and JavaScript, no build step; Leaflet for the map & Anyone can read and change them \\
    Delivery & Docker (\texttt{python:3.12-slim}), Docker Compose, GitHub Actions & One command to run; every pull
      request runs the tests \\
    Tests & pytest, httpx & 291 tests, including a race of twenty buyers for the last unit \\
    \bottomrule
  \end{tabular}
  \caption{Tech stack of the prototype.}
  \label{tab:stack}
\end{table}

The main swaps against the HLD's stack are SQLite for PostgreSQL, the outbox for Kafka and an in-process cache for
Redis (Table~\ref{tab:hld-vs-prototype}). Each sits behind an interface the full component can take over (the ORM, a
relay sink, a cache with get, set and delete), so swapping one should not need changes inside the services.

\section{Monitoring (HLD 12)}
\label{sec:m4-monitoring}

\texttt{GET /metrics} is a Prometheus scrape target with eight metrics: request time per route (a histogram, labelled
with the route template so the series do not grow with the number of stores), checkout holds by result, events
published by type, scans by result, cache lookups by hit or miss, holds in progress, the outbox lag and each store's
time since its last event. The app logs one JSON object per line with the request id. With
\texttt{docker compose --profile monitoring up}, Prometheus scrapes the app every 5\,s and evaluates the HLD 12
alert rules (Table~\ref{tab:alerts}), and Grafana opens on a dashboard of the same numbers. The ops page shows the
store view: problems first (stalled orders, bags waiting too long for a rider, items out of stock), then eleven KPI
cards, each beside its target and marked met or not met.

\begin{table}[ht]
  \centering
  \footnotesize
  \begin{tabular}{@{}>{\raggedright\arraybackslash}p{0.27\linewidth}>{\raggedright\arraybackslash}p{0.68\linewidth}@{}}
    \toprule
    Alert & Fires when \\
    \midrule
    Failed holds spike & More than 20\% of checkout holds fail over 5 minutes \\
    Event lag & The oldest unsent event is older than 30\,s for a minute \\
    Slow reads & The p99 of any read route is above the 100\,ms target \\
    Store offline & A store has published nothing for 15 minutes \\
    Cache and database disagree & Waits for the stock audit (task 1.8); not wired up \\
    \bottomrule
  \end{tabular}
  \caption{HLD 12 alert rules, in \texttt{monitoring/alerts.yml}. The thresholds are first guesses to tune with real
    traffic.}
  \label{tab:alerts}
\end{table}

\paragraph{Load test.} \texttt{scripts/load\_test.py} sends the HLD's mix: about 90\% reads and 10\% checkout holds,
a fifth of them paid and the rest released. Eight clients ran for 60\,s against one laptop, with the rate limit off
(Table~\ref{tab:load}). Reads met the 100\,ms target with room to spare. No call failed: the 189 refused holds are
the store correctly running out of stock. The outbox lag stayed at 0.2\,s, one relay pass.

\begin{table}[ht]
  \centering
  \small
  \begin{tabular}{@{}lrrrr@{}}
    \toprule
    Endpoint & Calls & Per second & p50 (ms) & p99 (ms) \\
    \midrule
    Availability (read)   & 8{,}494 & 141.6 & 18.2 &  89.9 \\
    Serviceability (read) & 8{,}640 & 144.0 & 10.7 &  26.1 \\
    Hold                  & 1{,}918 &  32.0 & 22.0 & 142.9 \\
    Payment               &     313 &   5.2 & 36.5 & 121.7 \\
    Release               & 1{,}416 &  23.6 & 18.1 & 142.8 \\
    \bottomrule
  \end{tabular}
  \caption{Load test: 8 clients for 60\,s, SQLite on one laptop, main at b39692e. No 5xx; 189 holds refused for
    lack of stock.}
  \label{tab:load}
\end{table}

What the numbers say, and what they do not:
\begin{itemize}
  \item The cache served 51\% of availability reads. This mix sends one stock write for about every two availability
    reads, and each write clears the store's entry; Member 1's benchmark at the HLD's 13 reads per write measured
    92\%.
  \item Writes are slower than reads because SQLite lets one writer in at a time. Under an earlier load test the
    picker's \texttt{next-task}, which claims an order and so writes, reached a p99 of 250--500\,ms. PostgreSQL's row
    locks would let writes to different stores proceed side by side.
  \item This is one laptop, not the HLD's cluster: the numbers show where the time goes, not what production would
    do.
\end{itemize}

\section{Trade-off: a transactional outbox instead of Kafka (HLD 13)}
\label{sec:m4-tradeoff}

The HLD sends events through Kafka. The prototype writes each event to a table in the same transaction as the change
that caused it, and a relay inside the app delivers it.

\textbf{What it gains.} An event exists if and only if its change committed: no event announces a hold that was
rolled back, and no committed sale goes unannounced. That is the hard part of using Kafka well, and the outbox gets
it from one database transaction. It needs no broker to run or operate, and under load the lag stayed at one relay
pass (0.2\,s).

\textbf{What it costs.}
\begin{itemize}
  \item \textbf{At least once, not exactly once.} A sink that fails gets the batch again, so consumers must ignore an
    event id they have seen. The edge-node experiment checks both cases (\texttt{tests/test\_edge\_node.py}). With
    the link down, 7 events queued; once it was back, all 7 arrived in order, once each. A crash after sending
    delivered a batch twice, and the receiver kept each event once.
  \item \textbf{One relay, inside the app.} If the app is down, nothing is relayed; but nothing changes either, and
    the backlog drains on restart.
  \item \textbf{One database.} The outbox shares SQLite's single writer, and \texttt{event\_log} is never pruned.
    Replay for a new consumer means reading the table, not rewinding a topic.
\end{itemize}

\textbf{When to switch.} A Kafka producer is one more sink (\texttt{events.add\_sink}), so the outbox would stay as
the reliable hand-off and Kafka would carry events to other services. That is worth it once several services
consume events or the stores move to edge nodes.

\section{Limits and open items}

\begin{table}[!htb]
  \centering
  \footnotesize
  \begin{tabular}{@{}>{\raggedright\arraybackslash}p{0.27\linewidth}>{\raggedright\arraybackslash}p{0.36\linewidth}
                  >{\raggedright\arraybackslash}p{0.31\linewidth}@{}}
    \toprule
    Limit & What it means & Next step \\
    \midrule
    No separate gateway, no TLS & The app does the gateway's work itself & A gateway in front, TLS everywhere \\
    \addlinespace[3pt]
    One shared signing key (HS256) & Whoever can check a token can also make one & Asymmetric signing, needed
      anyway for edge nodes \\
    \addlinespace[3pt]
    Tokens cannot be revoked & A stolen token works until it expires (30 minutes) & Short tokens with refresh, or a
      deny list \\
    \addlinespace[3pt]
    Rate limit per process and per address & Several instances or a proxy share or split the buckets & Move it to
      the gateway or Redis \\
    \addlinespace[3pt]
    Audit log never pruned & It grows with every staff call & Keep a set period, then archive \\
    \addlinespace[3pt]
    Zones rebuilt by hand; one store per cell & A busy store's outer customers are not passed to a neighbour & Let a
      cell list a second store \\
    \addlinespace[3pt]
    Demand is generated & The map and the site suggestions run on scripted orders & Real order history \\
    \addlinespace[3pt]
    No Alertmanager & Alerts show only on Prometheus's page & Route them to the on-call channel \\
    \bottomrule
  \end{tabular}
  \caption{Limits of the platform pieces and what would close them.}
  \label{tab:m4-limits}
\end{table}
. Needs: array, booktabs.
% member4-standalone.tex compiles these sections on their own. Numbers measured on main (b39692e), 21 November 2026.

\section{Architecture (HLD 4)}

The prototype is a \emph{modular monolith}: one FastAPI process that runs all eleven HLD services as eleven Python
packages under \texttt{app/modules/}. HLD 13.1 recommends this for a student build, and it keeps the HLD's service
boundaries, so each package can later move to its own process without a redesign. Every package has the same two
files: \texttt{service.py} holds the business logic as plain functions whose first argument is the database session,
and \texttt{router.py} holds the HTTP endpoints, which parse the input, call one service function and shape the
answer. Because the logic never touches HTTP, modules call each other directly and tests can call the logic without
a server.

Three rules hold the system together:
\begin{itemize}
  \item \textbf{One owner for stock.} Only the inventory module changes stock quantities: what is on hand, what is
    held and which bin it is in. Other modules ask it to hold, release, sell, receive or write off stock; forecasting
    and replenishment only set an item's reorder settings.
  \item \textbf{Hold at checkout, sell at packing.} Checkout raises \texttt{qty\_reserved}; nothing leaves
    \texttt{qty\_on\_hand} until the order is packed. A failed payment or a cancel only gives a hold back, so no stock
    count ever has to be ``put back''.
  \item \textbf{Check and change in one statement.} Every contested change is one conditional \texttt{UPDATE} with
    the check in its \texttt{WHERE} clause, so two requests can never both pass a check that only one of them should.
    The checkout hold, for example:
    \begin{center}\small
      \verb|UPDATE stock_levels SET qty_reserved = qty_reserved + :qty|\\
      \verb|WHERE store_id = :store AND sku_id = :sku AND qty_on_hand - qty_reserved >= :qty|
    \end{center}
\end{itemize}

Table~\ref{tab:hld-vs-prototype} lists where the prototype stands in for an HLD component.

\begin{table}[ht]
  \centering
  \footnotesize
  \begin{tabular}{@{}>{\raggedright\arraybackslash}p{0.2\linewidth}>{\raggedright\arraybackslash}p{0.44\linewidth}
                  >{\raggedright\arraybackslash}p{0.3\linewidth}@{}}
    \toprule
    HLD component & In the prototype & What a full build adds \\
    \midrule
    API gateway & The FastAPI app itself: login, roles, own-store checks, audit log, rate limit & A separate gateway
      with TLS and mutual TLS between services \\
    \addlinespace[3pt]
    PostgreSQL, sharded by store & One SQLite file (WAL mode), all stores in one database & PostgreSQL and
      migrations (task 1.7); sharding by \texttt{store\_id} \\
    \addlinespace[3pt]
    Redis ATP cache & A 2-second in-process cache of each store's availability list, cleared on every stock write
      & A shared Redis, so several app instances see one copy \\
    \addlinespace[3pt]
    Kafka & A transactional outbox (\texttt{event\_log}) and a relay inside the app that feeds each store's live feed
      & A Kafka producer as one more relay sink (section~\ref{sec:m4-tradeoff}) \\
    \addlinespace[3pt]
    Store edge node & A design note and an outage experiment (\texttt{docs/EDGE\_NODE.md}) & A second copy of the
      app in each store, syncing through the same relay \\
    \addlinespace[3pt]
    Prometheus, Grafana & Both, in Docker Compose, with the HLD 12 alert rules & An Alertmanager to send the
      alerts somewhere \\
    \bottomrule
  \end{tabular}
  \caption{HLD components and their stand-ins in the prototype.}
  \label{tab:hld-vs-prototype}
\end{table}

\section{Serviceability (HLD 5, service 1)}

Serviceability answers ``which store delivers here, and how fast?'' for every address a customer types, so it must
be cheap. Delivery zones are H3 hexagon cells at resolution 9: cells with edges of about 0.2\,km, so a zone border is
accurate to about 200\,m and a 3\,km zone is about 270 cells.
\begin{itemize}
  \item \textbf{Zones are built once.} \texttt{rebuild\_zones} gives each store every cell whose centre lies within its
    radius; where two stores claim a cell, the nearer one wins. The result is one row per cell.
  \item \textbf{A lookup is one row.} \texttt{check} turns the customer's point into its cell and reads that cell's
    row, instead of measuring the distance to every store.
  \item \textbf{Load shapes the answer.} A closed store serves nobody. A busy store (more than 5 unpacked orders)
    serves only the inner 70\% of its radius and quotes 5 minutes more. The ETA is 2 minutes to pack plus 3 minutes
    per kilometre, rounded up.
  \item \textbf{A miss is recorded.} A point no store serves is saved as its cell and the time, nothing else. These
    misses are the demand that network planning works from.
\end{itemize}
Under load (section~\ref{sec:m4-monitoring}) a serviceability check took 10.7\,ms at the median and 26.1\,ms at p99.

\section{Network Planning (HLD 5, service 10)}

Network Planning answers ``where should the next dark store go?'' (FR-12). It counts demand per H3 cell, never as a
name, phone number or exact address: a paid order carries the cell of its delivery address, and a turned-away
lookup is a miss. \texttt{GET /v1/network-planning/heatmap} returns both counts per cell over the last 30 days, and
the map at \texttt{/network.html} draws them over the current zones.

\texttt{suggest\_sites} chooses up to $k$ new sites by \emph{greedy maximum coverage} of the unmet demand. The
candidates are the cells with misses, except cells inside an existing zone and the cells listed in
\texttt{NO\_GO\_CELLS} (standing in for the zoning check of HLD 13.2). Each round picks the candidate whose 3\,km
radius covers the most unmet demand, breaking ties towards the more central site, then removes that demand and
repeats. It stops early when no candidate would gain anything. Greedy maximum coverage is guaranteed to cover at
least $1 - 1/e \approx 63\%$ of what the best choice of $k$ sites from the same candidates would cover. Only the 300
busiest miss cells are tried as sites, which bounds the work however many misses pile up.

\section{API and events (HLD 8)}

The API has 63 operations on 61 paths, all versioned under \texttt{/v1/} except \texttt{/health}, and documented by
OpenAPI at \texttt{/docs}.
Table~\ref{tab:api} groups them by service. They follow four conventions:
\begin{itemize}
  \item \textbf{One error shape.} Every refusal is \verb|{"error": {"code", "message", "details"}}|, with 409 when
    the request is valid but the current state forbids it (not enough stock, wrong job state), 422 for bad input and
    404 for something that does not exist. The code is stable, so a screen can react to it.
  \item \textbf{Order steps are safe to repeat.} Paying, packing, cancelling and handing over check the current state
    in the same conditional update, so a repeated call reports the state instead of acting twice. Stock adjustments
    take an \texttt{Idempotency-Key} header. A repeated hold or scan is not deduplicated: it holds or scans again.
  \item \textbf{Staff calls carry a token}; the customer calls (serviceability, availability, checkout) are open but
    rate-limited (section~\ref{sec:m4-security}).
  \item \textbf{Every response has a request id} (\texttt{X-Request-ID}), which also appears on every log line made
    while the request was handled.
\end{itemize}

\begin{table}[ht]
  \centering
  \small
  \begin{tabular}{@{}lrl@{}}
    \toprule
    Service & Operations & Examples \\
    \midrule
    Inventory           & 19 & availability, holds, adjustments, cycle counts, near-expiry \\
    Fulfilment          &  8 & cancel, handover, rider check-in, metrics, attention list \\
    Pick \& Pack        &  5 & next order or batch, scan, empty shelf, pack \\
    Replenishment       &  5 & run, dispatch, reorder settings \\
    Assortment          &  5 & catalogue, recommendations \\
    Forecasting         &  4 & run, compare models \\
    Dashboard           &  4 & store dashboard, city overview, KPIs, events \\
    Serviceability      &  3 & check, stores, rebuild zones \\
    Network planning    &  3 & heatmap, suggestions, coverage \\
    Admin               &  3 & reset demo data, audit log, health check \\
    Slotting, inbound, login & 4 & bins, suggest a bin, goods receipt, login \\
    \bottomrule
  \end{tabular}
  \caption{API operations by service.}
  \label{tab:api}
\end{table}

\paragraph{Events.} The code publishes 17 event types, such as \texttt{order.paid} when a customer pays and
\texttt{rider.assigned} when a bag is packed. Publishing writes the event to the \texttt{event\_log} table
\emph{inside the caller's transaction} and calls the subscribed handlers in that same transaction: this is
the transactional outbox of HLD 9.2. If the stock change rolls back, the event never existed;
if it commits, the event is stored. A relay inside the app reads unsent events every 0.2\,s, oldest first, up to 100
at a time, and hands each one to every sink. Today the only sink is the live feed: a WebSocket per store, so every
screen of that store sees an order paid or a line dropped within a second, without polling. An event is marked sent
only after every sink took it, so delivery is \emph{at least once}, and a consumer must ignore an event id it has
already seen.

\section{Security (HLD 10)}
\label{sec:m4-security}

\begin{itemize}
  \item \textbf{Passwords} are stored as salted scrypt hashes and compared in constant time; none is stored as typed.
  \item \textbf{Login} returns a signed token (JWT, HS256) that lasts 30 minutes and names the user, their role and
    their store. Six roles exist: picker, packer, receiver, manager, ops and admin.
  \item \textbf{Every staff endpoint} checks the role and, for store data, that the store is the caller's own:
    a picker of store 1 gets 403 for store 2. Ops may read every store, and \texttt{/v1/admin/} needs admin.
  \item \textbf{An audit log} keeps every call a logged-in user makes to a protected endpoint, refused or not: who,
    what (method, route and path), the status they got and the request id. An admin reads it at
    \texttt{GET /v1/audit}.
  \item \textbf{The open customer endpoints are rate-limited} per client address by a token bucket: 20 calls a second
    with bursts of 40, then 429 with a \texttt{Retry-After} header. This also limits how fast a caller can flood
    network planning with made-up misses.
  \item \textbf{Privacy.} Demand is kept as H3 cells, about 0.1\,km$^2$ each, so the planning data names no customer
    and no exact address. \texttt{/metrics} can be closed with \texttt{METRICS\_TOKEN}.
\end{itemize}
The limits are listed in Table~\ref{tab:m4-limits}: the built-in signing key and the shared demo password are public
for local development (the app warns about the key at start), a token cannot be revoked before it expires, and the
rate limit is
counted per process.

\section{Tech stack (HLD 11)}

\begin{table}[ht]
  \centering
  \footnotesize
  \begin{tabular}{@{}>{\raggedright\arraybackslash}p{0.22\linewidth}>{\raggedright\arraybackslash}p{0.4\linewidth}
                  >{\raggedright\arraybackslash}p{0.32\linewidth}@{}}
    \toprule
    Layer & Used & Why \\
    \midrule
    Language, web & Python 3.12, FastAPI, Uvicorn, Pydantic & Async server, validation and OpenAPI docs for free \\
    Data & SQLAlchemy 2 on SQLite (WAL) & Runs anywhere in one file; the same ORM code moves to PostgreSQL \\
    Geography & H3 (\texttt{h3} 4) & Equal-area hexagons make zones and demand one indexed lookup \\
    Login & PyJWT, scrypt from the standard library & No extra service for a prototype \\
    Monitoring & \texttt{prometheus-client}, Prometheus 2.53, Grafana 11.1 & The HLD's choice; one Compose profile \\
    Screens & HTML, CSS and JavaScript, no build step; Leaflet for the map & Anyone can read and change them \\
    Delivery & Docker (\texttt{python:3.12-slim}), Docker Compose, GitHub Actions & One command to run; every pull
      request runs the tests \\
    Tests & pytest, httpx & 291 tests, including a race of twenty buyers for the last unit \\
    \bottomrule
  \end{tabular}
  \caption{Tech stack of the prototype.}
  \label{tab:stack}
\end{table}

The main swaps against the HLD's stack are SQLite for PostgreSQL, the outbox for Kafka and an in-process cache for
Redis (Table~\ref{tab:hld-vs-prototype}). Each sits behind an interface the full component can take over (the ORM, a
relay sink, a cache with get, set and delete), so swapping one should not need changes inside the services.

\section{Monitoring (HLD 12)}
\label{sec:m4-monitoring}

\texttt{GET /metrics} is a Prometheus scrape target with eight metrics: request time per route (a histogram, labelled
with the route template so the series do not grow with the number of stores), checkout holds by result, events
published by type, scans by result, cache lookups by hit or miss, holds in progress, the outbox lag and each store's
time since its last event. The app logs one JSON object per line with the request id. With
\texttt{docker compose --profile monitoring up}, Prometheus scrapes the app every 5\,s and evaluates the HLD 12
alert rules (Table~\ref{tab:alerts}), and Grafana opens on a dashboard of the same numbers. The ops page shows the
store view: problems first (stalled orders, bags waiting too long for a rider, items out of stock), then eleven KPI
cards, each beside its target and marked met or not met.

\begin{table}[ht]
  \centering
  \footnotesize
  \begin{tabular}{@{}>{\raggedright\arraybackslash}p{0.27\linewidth}>{\raggedright\arraybackslash}p{0.68\linewidth}@{}}
    \toprule
    Alert & Fires when \\
    \midrule
    Failed holds spike & More than 20\% of checkout holds fail over 5 minutes \\
    Event lag & The oldest unsent event is older than 30\,s for a minute \\
    Slow reads & The p99 of any read route is above the 100\,ms target \\
    Store offline & A store has published nothing for 15 minutes \\
    Cache and database disagree & Waits for the stock audit (task 1.8); not wired up \\
    \bottomrule
  \end{tabular}
  \caption{HLD 12 alert rules, in \texttt{monitoring/alerts.yml}. The thresholds are first guesses to tune with real
    traffic.}
  \label{tab:alerts}
\end{table}

\paragraph{Load test.} \texttt{scripts/load\_test.py} sends the HLD's mix: about 90\% reads and 10\% checkout holds,
a fifth of them paid and the rest released. Eight clients ran for 60\,s against one laptop, with the rate limit off
(Table~\ref{tab:load}). Reads met the 100\,ms target with room to spare. No call failed: the 189 refused holds are
the store correctly running out of stock. The outbox lag stayed at 0.2\,s, one relay pass.

\begin{table}[ht]
  \centering
  \small
  \begin{tabular}{@{}lrrrr@{}}
    \toprule
    Endpoint & Calls & Per second & p50 (ms) & p99 (ms) \\
    \midrule
    Availability (read)   & 8{,}494 & 141.6 & 18.2 &  89.9 \\
    Serviceability (read) & 8{,}640 & 144.0 & 10.7 &  26.1 \\
    Hold                  & 1{,}918 &  32.0 & 22.0 & 142.9 \\
    Payment               &     313 &   5.2 & 36.5 & 121.7 \\
    Release               & 1{,}416 &  23.6 & 18.1 & 142.8 \\
    \bottomrule
  \end{tabular}
  \caption{Load test: 8 clients for 60\,s, SQLite on one laptop, main at b39692e. No 5xx; 189 holds refused for
    lack of stock.}
  \label{tab:load}
\end{table}

What the numbers say, and what they do not:
\begin{itemize}
  \item The cache served 51\% of availability reads. This mix sends one stock write for about every two availability
    reads, and each write clears the store's entry; Member 1's benchmark at the HLD's 13 reads per write measured
    92\%.
  \item Writes are slower than reads because SQLite lets one writer in at a time. Under an earlier load test the
    picker's \texttt{next-task}, which claims an order and so writes, reached a p99 of 250--500\,ms. PostgreSQL's row
    locks would let writes to different stores proceed side by side.
  \item This is one laptop, not the HLD's cluster: the numbers show where the time goes, not what production would
    do.
\end{itemize}

\section{Trade-off: a transactional outbox instead of Kafka (HLD 13)}
\label{sec:m4-tradeoff}

The HLD sends events through Kafka. The prototype writes each event to a table in the same transaction as the change
that caused it, and a relay inside the app delivers it.

\textbf{What it gains.} An event exists if and only if its change committed: no event announces a hold that was
rolled back, and no committed sale goes unannounced. That is the hard part of using Kafka well, and the outbox gets
it from one database transaction. It needs no broker to run or operate, and under load the lag stayed at one relay
pass (0.2\,s).

\textbf{What it costs.}
\begin{itemize}
  \item \textbf{At least once, not exactly once.} A sink that fails gets the batch again, so consumers must ignore an
    event id they have seen. The edge-node experiment checks both cases (\texttt{tests/test\_edge\_node.py}). With
    the link down, 7 events queued; once it was back, all 7 arrived in order, once each. A crash after sending
    delivered a batch twice, and the receiver kept each event once.
  \item \textbf{One relay, inside the app.} If the app is down, nothing is relayed; but nothing changes either, and
    the backlog drains on restart.
  \item \textbf{One database.} The outbox shares SQLite's single writer, and \texttt{event\_log} is never pruned.
    Replay for a new consumer means reading the table, not rewinding a topic.
\end{itemize}

\textbf{When to switch.} A Kafka producer is one more sink (\texttt{events.add\_sink}), so the outbox would stay as
the reliable hand-off and Kafka would carry events to other services. That is worth it once several services
consume events or the stores move to edge nodes.

\section{Limits and open items}

\begin{table}[!htb]
  \centering
  \footnotesize
  \begin{tabular}{@{}>{\raggedright\arraybackslash}p{0.27\linewidth}>{\raggedright\arraybackslash}p{0.36\linewidth}
                  >{\raggedright\arraybackslash}p{0.31\linewidth}@{}}
    \toprule
    Limit & What it means & Next step \\
    \midrule
    No separate gateway, no TLS & The app does the gateway's work itself & A gateway in front, TLS everywhere \\
    \addlinespace[3pt]
    One shared signing key (HS256) & Whoever can check a token can also make one & Asymmetric signing, needed
      anyway for edge nodes \\
    \addlinespace[3pt]
    Tokens cannot be revoked & A stolen token works until it expires (30 minutes) & Short tokens with refresh, or a
      deny list \\
    \addlinespace[3pt]
    Rate limit per process and per address & Several instances or a proxy share or split the buckets & Move it to
      the gateway or Redis \\
    \addlinespace[3pt]
    Audit log never pruned & It grows with every staff call & Keep a set period, then archive \\
    \addlinespace[3pt]
    Zones rebuilt by hand; one store per cell & A busy store's outer customers are not passed to a neighbour & Let a
      cell list a second store \\
    \addlinespace[3pt]
    Demand is generated & The map and the site suggestions run on scripted orders & Real order history \\
    \addlinespace[3pt]
    No Alertmanager & Alerts show only on Prometheus's page & Route them to the on-call channel \\
    \bottomrule
  \end{tabular}
  \caption{Limits of the platform pieces and what would close them.}
  \label{tab:m4-limits}
\end{table}
. Needs: array, booktabs.
% member4-standalone.tex compiles these sections on their own. Numbers measured on main (b39692e), 21 November 2026.

\section{Architecture (HLD 4)}

The prototype is a \emph{modular monolith}: one FastAPI process that runs all eleven HLD services as eleven Python
packages under \texttt{app/modules/}. HLD 13.1 recommends this for a student build, and it keeps the HLD's service
boundaries, so each package can later move to its own process without a redesign. Every package has the same two
files: \texttt{service.py} holds the business logic as plain functions whose first argument is the database session,
and \texttt{router.py} holds the HTTP endpoints, which parse the input, call one service function and shape the
answer. Because the logic never touches HTTP, modules call each other directly and tests can call the logic without
a server.

Three rules hold the system together:
\begin{itemize}
  \item \textbf{One owner for stock.} Only the inventory module changes stock quantities: what is on hand, what is
    held and which bin it is in. Other modules ask it to hold, release, sell, receive or write off stock; forecasting
    and replenishment only set an item's reorder settings.
  \item \textbf{Hold at checkout, sell at packing.} Checkout raises \texttt{qty\_reserved}; nothing leaves
    \texttt{qty\_on\_hand} until the order is packed. A failed payment or a cancel only gives a hold back, so no stock
    count ever has to be ``put back''.
  \item \textbf{Check and change in one statement.} Every contested change is one conditional \texttt{UPDATE} with
    the check in its \texttt{WHERE} clause, so two requests can never both pass a check that only one of them should.
    The checkout hold, for example:
    \begin{center}\small
      \verb|UPDATE stock_levels SET qty_reserved = qty_reserved + :qty|\\
      \verb|WHERE store_id = :store AND sku_id = :sku AND qty_on_hand - qty_reserved >= :qty|
    \end{center}
\end{itemize}

Table~\ref{tab:hld-vs-prototype} lists where the prototype stands in for an HLD component.

\begin{table}[ht]
  \centering
  \footnotesize
  \begin{tabular}{@{}>{\raggedright\arraybackslash}p{0.2\linewidth}>{\raggedright\arraybackslash}p{0.44\linewidth}
                  >{\raggedright\arraybackslash}p{0.3\linewidth}@{}}
    \toprule
    HLD component & In the prototype & What a full build adds \\
    \midrule
    API gateway & The FastAPI app itself: login, roles, own-store checks, audit log, rate limit & A separate gateway
      with TLS and mutual TLS between services \\
    \addlinespace[3pt]
    PostgreSQL, sharded by store & One SQLite file (WAL mode), all stores in one database & PostgreSQL and
      migrations (task 1.7); sharding by \texttt{store\_id} \\
    \addlinespace[3pt]
    Redis ATP cache & A 2-second in-process cache of each store's availability list, cleared on every stock write
      & A shared Redis, so several app instances see one copy \\
    \addlinespace[3pt]
    Kafka & A transactional outbox (\texttt{event\_log}) and a relay inside the app that feeds each store's live feed
      & A Kafka producer as one more relay sink (section~\ref{sec:m4-tradeoff}) \\
    \addlinespace[3pt]
    Store edge node & A design note and an outage experiment (\texttt{docs/EDGE\_NODE.md}) & A second copy of the
      app in each store, syncing through the same relay \\
    \addlinespace[3pt]
    Prometheus, Grafana & Both, in Docker Compose, with the HLD 12 alert rules & An Alertmanager to send the
      alerts somewhere \\
    \bottomrule
  \end{tabular}
  \caption{HLD components and their stand-ins in the prototype.}
  \label{tab:hld-vs-prototype}
\end{table}

\section{Serviceability (HLD 5, service 1)}

Serviceability answers ``which store delivers here, and how fast?'' for every address a customer types, so it must
be cheap. Delivery zones are H3 hexagon cells at resolution 9: cells with edges of about 0.2\,km, so a zone border is
accurate to about 200\,m and a 3\,km zone is about 270 cells.
\begin{itemize}
  \item \textbf{Zones are built once.} \texttt{rebuild\_zones} gives each store every cell whose centre lies within its
    radius; where two stores claim a cell, the nearer one wins. The result is one row per cell.
  \item \textbf{A lookup is one row.} \texttt{check} turns the customer's point into its cell and reads that cell's
    row, instead of measuring the distance to every store.
  \item \textbf{Load shapes the answer.} A closed store serves nobody. A busy store (more than 5 unpacked orders)
    serves only the inner 70\% of its radius and quotes 5 minutes more. The ETA is 2 minutes to pack plus 3 minutes
    per kilometre, rounded up.
  \item \textbf{A miss is recorded.} A point no store serves is saved as its cell and the time, nothing else. These
    misses are the demand that network planning works from.
\end{itemize}
Under load (section~\ref{sec:m4-monitoring}) a serviceability check took 10.7\,ms at the median and 26.1\,ms at p99.

\section{Network Planning (HLD 5, service 10)}

Network Planning answers ``where should the next dark store go?'' (FR-12). It counts demand per H3 cell, never as a
name, phone number or exact address: a paid order carries the cell of its delivery address, and a turned-away
lookup is a miss. \texttt{GET /v1/network-planning/heatmap} returns both counts per cell over the last 30 days, and
the map at \texttt{/network.html} draws them over the current zones.

\texttt{suggest\_sites} chooses up to $k$ new sites by \emph{greedy maximum coverage} of the unmet demand. The
candidates are the cells with misses, except cells inside an existing zone and the cells listed in
\texttt{NO\_GO\_CELLS} (standing in for the zoning check of HLD 13.2). Each round picks the candidate whose 3\,km
radius covers the most unmet demand, breaking ties towards the more central site, then removes that demand and
repeats. It stops early when no candidate would gain anything. Greedy maximum coverage is guaranteed to cover at
least $1 - 1/e \approx 63\%$ of what the best choice of $k$ sites from the same candidates would cover. Only the 300
busiest miss cells are tried as sites, which bounds the work however many misses pile up.

\section{API and events (HLD 8)}

The API has 63 operations on 61 paths, all versioned under \texttt{/v1/} except \texttt{/health}, and documented by
OpenAPI at \texttt{/docs}.
Table~\ref{tab:api} groups them by service. They follow four conventions:
\begin{itemize}
  \item \textbf{One error shape.} Every refusal is \verb|{"error": {"code", "message", "details"}}|, with 409 when
    the request is valid but the current state forbids it (not enough stock, wrong job state), 422 for bad input and
    404 for something that does not exist. The code is stable, so a screen can react to it.
  \item \textbf{Order steps are safe to repeat.} Paying, packing, cancelling and handing over check the current state
    in the same conditional update, so a repeated call reports the state instead of acting twice. Stock adjustments
    take an \texttt{Idempotency-Key} header. A repeated hold or scan is not deduplicated: it holds or scans again.
  \item \textbf{Staff calls carry a token}; the customer calls (serviceability, availability, checkout) are open but
    rate-limited (section~\ref{sec:m4-security}).
  \item \textbf{Every response has a request id} (\texttt{X-Request-ID}), which also appears on every log line made
    while the request was handled.
\end{itemize}

\begin{table}[ht]
  \centering
  \small
  \begin{tabular}{@{}lrl@{}}
    \toprule
    Service & Operations & Examples \\
    \midrule
    Inventory           & 19 & availability, holds, adjustments, cycle counts, near-expiry \\
    Fulfilment          &  8 & cancel, handover, rider check-in, metrics, attention list \\
    Pick \& Pack        &  5 & next order or batch, scan, empty shelf, pack \\
    Replenishment       &  5 & run, dispatch, reorder settings \\
    Assortment          &  5 & catalogue, recommendations \\
    Forecasting         &  4 & run, compare models \\
    Dashboard           &  4 & store dashboard, city overview, KPIs, events \\
    Serviceability      &  3 & check, stores, rebuild zones \\
    Network planning    &  3 & heatmap, suggestions, coverage \\
    Admin               &  3 & reset demo data, audit log, health check \\
    Slotting, inbound, login & 4 & bins, suggest a bin, goods receipt, login \\
    \bottomrule
  \end{tabular}
  \caption{API operations by service.}
  \label{tab:api}
\end{table}

\paragraph{Events.} The code publishes 17 event types, such as \texttt{order.paid} when a customer pays and
\texttt{rider.assigned} when a bag is packed. Publishing writes the event to the \texttt{event\_log} table
\emph{inside the caller's transaction} and calls the subscribed handlers in that same transaction: this is
the transactional outbox of HLD 9.2. If the stock change rolls back, the event never existed;
if it commits, the event is stored. A relay inside the app reads unsent events every 0.2\,s, oldest first, up to 100
at a time, and hands each one to every sink. Today the only sink is the live feed: a WebSocket per store, so every
screen of that store sees an order paid or a line dropped within a second, without polling. An event is marked sent
only after every sink took it, so delivery is \emph{at least once}, and a consumer must ignore an event id it has
already seen.

\section{Security (HLD 10)}
\label{sec:m4-security}

\begin{itemize}
  \item \textbf{Passwords} are stored as salted scrypt hashes and compared in constant time; none is stored as typed.
  \item \textbf{Login} returns a signed token (JWT, HS256) that lasts 30 minutes and names the user, their role and
    their store. Six roles exist: picker, packer, receiver, manager, ops and admin.
  \item \textbf{Every staff endpoint} checks the role and, for store data, that the store is the caller's own:
    a picker of store 1 gets 403 for store 2. Ops may read every store, and \texttt{/v1/admin/} needs admin.
  \item \textbf{An audit log} keeps every call a logged-in user makes to a protected endpoint, refused or not: who,
    what (method, route and path), the status they got and the request id. An admin reads it at
    \texttt{GET /v1/audit}.
  \item \textbf{The open customer endpoints are rate-limited} per client address by a token bucket: 20 calls a second
    with bursts of 40, then 429 with a \texttt{Retry-After} header. This also limits how fast a caller can flood
    network planning with made-up misses.
  \item \textbf{Privacy.} Demand is kept as H3 cells, about 0.1\,km$^2$ each, so the planning data names no customer
    and no exact address. \texttt{/metrics} can be closed with \texttt{METRICS\_TOKEN}.
\end{itemize}
The limits are listed in Table~\ref{tab:m4-limits}: the built-in signing key and the shared demo password are public
for local development (the app warns about the key at start), a token cannot be revoked before it expires, and the
rate limit is
counted per process.

\section{Tech stack (HLD 11)}

\begin{table}[ht]
  \centering
  \footnotesize
  \begin{tabular}{@{}>{\raggedright\arraybackslash}p{0.22\linewidth}>{\raggedright\arraybackslash}p{0.4\linewidth}
                  >{\raggedright\arraybackslash}p{0.32\linewidth}@{}}
    \toprule
    Layer & Used & Why \\
    \midrule
    Language, web & Python 3.12, FastAPI, Uvicorn, Pydantic & Async server, validation and OpenAPI docs for free \\
    Data & SQLAlchemy 2 on SQLite (WAL) & Runs anywhere in one file; the same ORM code moves to PostgreSQL \\
    Geography & H3 (\texttt{h3} 4) & Equal-area hexagons make zones and demand one indexed lookup \\
    Login & PyJWT, scrypt from the standard library & No extra service for a prototype \\
    Monitoring & \texttt{prometheus-client}, Prometheus 2.53, Grafana 11.1 & The HLD's choice; one Compose profile \\
    Screens & HTML, CSS and JavaScript, no build step; Leaflet for the map & Anyone can read and change them \\
    Delivery & Docker (\texttt{python:3.12-slim}), Docker Compose, GitHub Actions & One command to run; every pull
      request runs the tests \\
    Tests & pytest, httpx & 291 tests, including a race of twenty buyers for the last unit \\
    \bottomrule
  \end{tabular}
  \caption{Tech stack of the prototype.}
  \label{tab:stack}
\end{table}

The main swaps against the HLD's stack are SQLite for PostgreSQL, the outbox for Kafka and an in-process cache for
Redis (Table~\ref{tab:hld-vs-prototype}). Each sits behind an interface the full component can take over (the ORM, a
relay sink, a cache with get, set and delete), so swapping one should not need changes inside the services.

\section{Monitoring (HLD 12)}
\label{sec:m4-monitoring}

\texttt{GET /metrics} is a Prometheus scrape target with eight metrics: request time per route (a histogram, labelled
with the route template so the series do not grow with the number of stores), checkout holds by result, events
published by type, scans by result, cache lookups by hit or miss, holds in progress, the outbox lag and each store's
time since its last event. The app logs one JSON object per line with the request id. With
\texttt{docker compose --profile monitoring up}, Prometheus scrapes the app every 5\,s and evaluates the HLD 12
alert rules (Table~\ref{tab:alerts}), and Grafana opens on a dashboard of the same numbers. The ops page shows the
store view: problems first (stalled orders, bags waiting too long for a rider, items out of stock), then eleven KPI
cards, each beside its target and marked met or not met.

\begin{table}[ht]
  \centering
  \footnotesize
  \begin{tabular}{@{}>{\raggedright\arraybackslash}p{0.27\linewidth}>{\raggedright\arraybackslash}p{0.68\linewidth}@{}}
    \toprule
    Alert & Fires when \\
    \midrule
    Failed holds spike & More than 20\% of checkout holds fail over 5 minutes \\
    Event lag & The oldest unsent event is older than 30\,s for a minute \\
    Slow reads & The p99 of any read route is above the 100\,ms target \\
    Store offline & A store has published nothing for 15 minutes \\
    Cache and database disagree & Waits for the stock audit (task 1.8); not wired up \\
    \bottomrule
  \end{tabular}
  \caption{HLD 12 alert rules, in \texttt{monitoring/alerts.yml}. The thresholds are first guesses to tune with real
    traffic.}
  \label{tab:alerts}
\end{table}

\paragraph{Load test.} \texttt{scripts/load\_test.py} sends the HLD's mix: about 90\% reads and 10\% checkout holds,
a fifth of them paid and the rest released. Eight clients ran for 60\,s against one laptop, with the rate limit off
(Table~\ref{tab:load}). Reads met the 100\,ms target with room to spare. No call failed: the 189 refused holds are
the store correctly running out of stock. The outbox lag stayed at 0.2\,s, one relay pass.

\begin{table}[ht]
  \centering
  \small
  \begin{tabular}{@{}lrrrr@{}}
    \toprule
    Endpoint & Calls & Per second & p50 (ms) & p99 (ms) \\
    \midrule
    Availability (read)   & 8{,}494 & 141.6 & 18.2 &  89.9 \\
    Serviceability (read) & 8{,}640 & 144.0 & 10.7 &  26.1 \\
    Hold                  & 1{,}918 &  32.0 & 22.0 & 142.9 \\
    Payment               &     313 &   5.2 & 36.5 & 121.7 \\
    Release               & 1{,}416 &  23.6 & 18.1 & 142.8 \\
    \bottomrule
  \end{tabular}
  \caption{Load test: 8 clients for 60\,s, SQLite on one laptop, main at b39692e. No 5xx; 189 holds refused for
    lack of stock.}
  \label{tab:load}
\end{table}

What the numbers say, and what they do not:
\begin{itemize}
  \item The cache served 51\% of availability reads. This mix sends one stock write for about every two availability
    reads, and each write clears the store's entry; Member 1's benchmark at the HLD's 13 reads per write measured
    92\%.
  \item Writes are slower than reads because SQLite lets one writer in at a time. Under an earlier load test the
    picker's \texttt{next-task}, which claims an order and so writes, reached a p99 of 250--500\,ms. PostgreSQL's row
    locks would let writes to different stores proceed side by side.
  \item This is one laptop, not the HLD's cluster: the numbers show where the time goes, not what production would
    do.
\end{itemize}

\section{Trade-off: a transactional outbox instead of Kafka (HLD 13)}
\label{sec:m4-tradeoff}

The HLD sends events through Kafka. The prototype writes each event to a table in the same transaction as the change
that caused it, and a relay inside the app delivers it.

\textbf{What it gains.} An event exists if and only if its change committed: no event announces a hold that was
rolled back, and no committed sale goes unannounced. That is the hard part of using Kafka well, and the outbox gets
it from one database transaction. It needs no broker to run or operate, and under load the lag stayed at one relay
pass (0.2\,s).

\textbf{What it costs.}
\begin{itemize}
  \item \textbf{At least once, not exactly once.} A sink that fails gets the batch again, so consumers must ignore an
    event id they have seen. The edge-node experiment checks both cases (\texttt{tests/test\_edge\_node.py}). With
    the link down, 7 events queued; once it was back, all 7 arrived in order, once each. A crash after sending
    delivered a batch twice, and the receiver kept each event once.
  \item \textbf{One relay, inside the app.} If the app is down, nothing is relayed; but nothing changes either, and
    the backlog drains on restart.
  \item \textbf{One database.} The outbox shares SQLite's single writer, and \texttt{event\_log} is never pruned.
    Replay for a new consumer means reading the table, not rewinding a topic.
\end{itemize}

\textbf{When to switch.} A Kafka producer is one more sink (\texttt{events.add\_sink}), so the outbox would stay as
the reliable hand-off and Kafka would carry events to other services. That is worth it once several services
consume events or the stores move to edge nodes.

\section{Limits and open items}

\begin{table}[!htb]
  \centering
  \footnotesize
  \begin{tabular}{@{}>{\raggedright\arraybackslash}p{0.27\linewidth}>{\raggedright\arraybackslash}p{0.36\linewidth}
                  >{\raggedright\arraybackslash}p{0.31\linewidth}@{}}
    \toprule
    Limit & What it means & Next step \\
    \midrule
    No separate gateway, no TLS & The app does the gateway's work itself & A gateway in front, TLS everywhere \\
    \addlinespace[3pt]
    One shared signing key (HS256) & Whoever can check a token can also make one & Asymmetric signing, needed
      anyway for edge nodes \\
    \addlinespace[3pt]
    Tokens cannot be revoked & A stolen token works until it expires (30 minutes) & Short tokens with refresh, or a
      deny list \\
    \addlinespace[3pt]
    Rate limit per process and per address & Several instances or a proxy share or split the buckets & Move it to
      the gateway or Redis \\
    \addlinespace[3pt]
    Audit log never pruned & It grows with every staff call & Keep a set period, then archive \\
    \addlinespace[3pt]
    Zones rebuilt by hand; one store per cell & A busy store's outer customers are not passed to a neighbour & Let a
      cell list a second store \\
    \addlinespace[3pt]
    Demand is generated & The map and the site suggestions run on scripted orders & Real order history \\
    \addlinespace[3pt]
    No Alertmanager & Alerts show only on Prometheus's page & Route them to the on-call channel \\
    \bottomrule
  \end{tabular}
  \caption{Limits of the platform pieces and what would close them.}
  \label{tab:m4-limits}
\end{table}
. Needs: array, booktabs.
% member4-standalone.tex compiles these sections on their own. Numbers measured on main (b39692e), 21 November 2026.

\section{Architecture (HLD 4)}

The prototype is a \emph{modular monolith}: one FastAPI process that runs all eleven HLD services as eleven Python
packages under \texttt{app/modules/}. HLD 13.1 recommends this for a student build, and it keeps the HLD's service
boundaries, so each package can later move to its own process without a redesign. Every package has the same two
files: \texttt{service.py} holds the business logic as plain functions whose first argument is the database session,
and \texttt{router.py} holds the HTTP endpoints, which parse the input, call one service function and shape the
answer. Because the logic never touches HTTP, modules call each other directly and tests can call the logic without
a server.

Three rules hold the system together:
\begin{itemize}
  \item \textbf{One owner for stock.} Only the inventory module changes stock quantities: what is on hand, what is
    held and which bin it is in. Other modules ask it to hold, release, sell, receive or write off stock; forecasting
    and replenishment only set an item's reorder settings.
  \item \textbf{Hold at checkout, sell at packing.} Checkout raises \texttt{qty\_reserved}; nothing leaves
    \texttt{qty\_on\_hand} until the order is packed. A failed payment or a cancel only gives a hold back, so no stock
    count ever has to be ``put back''.
  \item \textbf{Check and change in one statement.} Every contested change is one conditional \texttt{UPDATE} with
    the check in its \texttt{WHERE} clause, so two requests can never both pass a check that only one of them should.
    The checkout hold, for example:
    \begin{center}\small
      \verb|UPDATE stock_levels SET qty_reserved = qty_reserved + :qty|\\
      \verb|WHERE store_id = :store AND sku_id = :sku AND qty_on_hand - qty_reserved >= :qty|
    \end{center}
\end{itemize}

Table~\ref{tab:hld-vs-prototype} lists where the prototype stands in for an HLD component.

\begin{table}[ht]
  \centering
  \footnotesize
  \begin{tabular}{@{}>{\raggedright\arraybackslash}p{0.2\linewidth}>{\raggedright\arraybackslash}p{0.44\linewidth}
                  >{\raggedright\arraybackslash}p{0.3\linewidth}@{}}
    \toprule
    HLD component & In the prototype & What a full build adds \\
    \midrule
    API gateway & The FastAPI app itself: login, roles, own-store checks, audit log, rate limit & A separate gateway
      with TLS and mutual TLS between services \\
    \addlinespace[3pt]
    PostgreSQL, sharded by store & One SQLite file (WAL mode), all stores in one database & PostgreSQL and
      migrations (task 1.7); sharding by \texttt{store\_id} \\
    \addlinespace[3pt]
    Redis ATP cache & A 2-second in-process cache of each store's availability list, cleared on every stock write
      & A shared Redis, so several app instances see one copy \\
    \addlinespace[3pt]
    Kafka & A transactional outbox (\texttt{event\_log}) and a relay inside the app that feeds each store's live feed
      & A Kafka producer as one more relay sink (section~\ref{sec:m4-tradeoff}) \\
    \addlinespace[3pt]
    Store edge node & A design note and an outage experiment (\texttt{docs/EDGE\_NODE.md}) & A second copy of the
      app in each store, syncing through the same relay \\
    \addlinespace[3pt]
    Prometheus, Grafana & Both, in Docker Compose, with the HLD 12 alert rules & An Alertmanager to send the
      alerts somewhere \\
    \bottomrule
  \end{tabular}
  \caption{HLD components and their stand-ins in the prototype.}
  \label{tab:hld-vs-prototype}
\end{table}

\section{Serviceability (HLD 5, service 1)}

Serviceability answers ``which store delivers here, and how fast?'' for every address a customer types, so it must
be cheap. Delivery zones are H3 hexagon cells at resolution 9: cells with edges of about 0.2\,km, so a zone border is
accurate to about 200\,m and a 3\,km zone is about 270 cells.
\begin{itemize}
  \item \textbf{Zones are built once.} \texttt{rebuild\_zones} gives each store every cell whose centre lies within its
    radius; where two stores claim a cell, the nearer one wins. The result is one row per cell.
  \item \textbf{A lookup is one row.} \texttt{check} turns the customer's point into its cell and reads that cell's
    row, instead of measuring the distance to every store.
  \item \textbf{Load shapes the answer.} A closed store serves nobody. A busy store (more than 5 unpacked orders)
    serves only the inner 70\% of its radius and quotes 5 minutes more. The ETA is 2 minutes to pack plus 3 minutes
    per kilometre, rounded up.
  \item \textbf{A miss is recorded.} A point no store serves is saved as its cell and the time, nothing else. These
    misses are the demand that network planning works from.
\end{itemize}
Under load (section~\ref{sec:m4-monitoring}) a serviceability check took 10.7\,ms at the median and 26.1\,ms at p99.

\section{Network Planning (HLD 5, service 10)}

Network Planning answers ``where should the next dark store go?'' (FR-12). It counts demand per H3 cell, never as a
name, phone number or exact address: a paid order carries the cell of its delivery address, and a turned-away
lookup is a miss. \texttt{GET /v1/network-planning/heatmap} returns both counts per cell over the last 30 days, and
the map at \texttt{/network.html} draws them over the current zones.

\texttt{suggest\_sites} chooses up to $k$ new sites by \emph{greedy maximum coverage} of the unmet demand. The
candidates are the cells with misses, except cells inside an existing zone and the cells listed in
\texttt{NO\_GO\_CELLS} (standing in for the zoning check of HLD 13.2). Each round picks the candidate whose 3\,km
radius covers the most unmet demand, breaking ties towards the more central site, then removes that demand and
repeats. It stops early when no candidate would gain anything. Greedy maximum coverage is guaranteed to cover at
least $1 - 1/e \approx 63\%$ of what the best choice of $k$ sites from the same candidates would cover. Only the 300
busiest miss cells are tried as sites, which bounds the work however many misses pile up.

\section{API and events (HLD 8)}

The API has 63 operations on 61 paths, all versioned under \texttt{/v1/} except \texttt{/health}, and documented by
OpenAPI at \texttt{/docs}.
Table~\ref{tab:api} groups them by service. They follow four conventions:
\begin{itemize}
  \item \textbf{One error shape.} Every refusal is \verb|{"error": {"code", "message", "details"}}|, with 409 when
    the request is valid but the current state forbids it (not enough stock, wrong job state), 422 for bad input and
    404 for something that does not exist. The code is stable, so a screen can react to it.
  \item \textbf{Order steps are safe to repeat.} Paying, packing, cancelling and handing over check the current state
    in the same conditional update, so a repeated call reports the state instead of acting twice. Stock adjustments
    take an \texttt{Idempotency-Key} header. A repeated hold or scan is not deduplicated: it holds or scans again.
  \item \textbf{Staff calls carry a token}; the customer calls (serviceability, availability, checkout) are open but
    rate-limited (section~\ref{sec:m4-security}).
  \item \textbf{Every response has a request id} (\texttt{X-Request-ID}), which also appears on every log line made
    while the request was handled.
\end{itemize}

\begin{table}[ht]
  \centering
  \small
  \begin{tabular}{@{}lrl@{}}
    \toprule
    Service & Operations & Examples \\
    \midrule
    Inventory           & 19 & availability, holds, adjustments, cycle counts, near-expiry \\
    Fulfilment          &  8 & cancel, handover, rider check-in, metrics, attention list \\
    Pick \& Pack        &  5 & next order or batch, scan, empty shelf, pack \\
    Replenishment       &  5 & run, dispatch, reorder settings \\
    Assortment          &  5 & catalogue, recommendations \\
    Forecasting         &  4 & run, compare models \\
    Dashboard           &  4 & store dashboard, city overview, KPIs, events \\
    Serviceability      &  3 & check, stores, rebuild zones \\
    Network planning    &  3 & heatmap, suggestions, coverage \\
    Admin               &  3 & reset demo data, audit log, health check \\
    Slotting, inbound, login & 4 & bins, suggest a bin, goods receipt, login \\
    \bottomrule
  \end{tabular}
  \caption{API operations by service.}
  \label{tab:api}
\end{table}

\paragraph{Events.} The code publishes 17 event types, such as \texttt{order.paid} when a customer pays and
\texttt{rider.assigned} when a bag is packed. Publishing writes the event to the \texttt{event\_log} table
\emph{inside the caller's transaction} and calls the subscribed handlers in that same transaction: this is
the transactional outbox of HLD 9.2. If the stock change rolls back, the event never existed;
if it commits, the event is stored. A relay inside the app reads unsent events every 0.2\,s, oldest first, up to 100
at a time, and hands each one to every sink. Today the only sink is the live feed: a WebSocket per store, so every
screen of that store sees an order paid or a line dropped within a second, without polling. An event is marked sent
only after every sink took it, so delivery is \emph{at least once}, and a consumer must ignore an event id it has
already seen.

\section{Security (HLD 10)}
\label{sec:m4-security}

\begin{itemize}
  \item \textbf{Passwords} are stored as salted scrypt hashes and compared in constant time; none is stored as typed.
  \item \textbf{Login} returns a signed token (JWT, HS256) that lasts 30 minutes and names the user, their role and
    their store. Six roles exist: picker, packer, receiver, manager, ops and admin.
  \item \textbf{Every staff endpoint} checks the role and, for store data, that the store is the caller's own:
    a picker of store 1 gets 403 for store 2. Ops may read every store, and \texttt{/v1/admin/} needs admin.
  \item \textbf{An audit log} keeps every call a logged-in user makes to a protected endpoint, refused or not: who,
    what (method, route and path), the status they got and the request id. An admin reads it at
    \texttt{GET /v1/audit}.
  \item \textbf{The open customer endpoints are rate-limited} per client address by a token bucket: 20 calls a second
    with bursts of 40, then 429 with a \texttt{Retry-After} header. This also limits how fast a caller can flood
    network planning with made-up misses.
  \item \textbf{Privacy.} Demand is kept as H3 cells, about 0.1\,km$^2$ each, so the planning data names no customer
    and no exact address. \texttt{/metrics} can be closed with \texttt{METRICS\_TOKEN}.
\end{itemize}
The limits are listed in Table~\ref{tab:m4-limits}: the built-in signing key and the shared demo password are public
for local development (the app warns about the key at start), a token cannot be revoked before it expires, and the
rate limit is
counted per process.

\section{Tech stack (HLD 11)}

\begin{table}[ht]
  \centering
  \footnotesize
  \begin{tabular}{@{}>{\raggedright\arraybackslash}p{0.22\linewidth}>{\raggedright\arraybackslash}p{0.4\linewidth}
                  >{\raggedright\arraybackslash}p{0.32\linewidth}@{}}
    \toprule
    Layer & Used & Why \\
    \midrule
    Language, web & Python 3.12, FastAPI, Uvicorn, Pydantic & Async server, validation and OpenAPI docs for free \\
    Data & SQLAlchemy 2 on SQLite (WAL) & Runs anywhere in one file; the same ORM code moves to PostgreSQL \\
    Geography & H3 (\texttt{h3} 4) & Equal-area hexagons make zones and demand one indexed lookup \\
    Login & PyJWT, scrypt from the standard library & No extra service for a prototype \\
    Monitoring & \texttt{prometheus-client}, Prometheus 2.53, Grafana 11.1 & The HLD's choice; one Compose profile \\
    Screens & HTML, CSS and JavaScript, no build step; Leaflet for the map & Anyone can read and change them \\
    Delivery & Docker (\texttt{python:3.12-slim}), Docker Compose, GitHub Actions & One command to run; every pull
      request runs the tests \\
    Tests & pytest, httpx & 291 tests, including a race of twenty buyers for the last unit \\
    \bottomrule
  \end{tabular}
  \caption{Tech stack of the prototype.}
  \label{tab:stack}
\end{table}

The main swaps against the HLD's stack are SQLite for PostgreSQL, the outbox for Kafka and an in-process cache for
Redis (Table~\ref{tab:hld-vs-prototype}). Each sits behind an interface the full component can take over (the ORM, a
relay sink, a cache with get, set and delete), so swapping one should not need changes inside the services.

\section{Monitoring (HLD 12)}
\label{sec:m4-monitoring}

\texttt{GET /metrics} is a Prometheus scrape target with eight metrics: request time per route (a histogram, labelled
with the route template so the series do not grow with the number of stores), checkout holds by result, events
published by type, scans by result, cache lookups by hit or miss, holds in progress, the outbox lag and each store's
time since its last event. The app logs one JSON object per line with the request id. With
\texttt{docker compose --profile monitoring up}, Prometheus scrapes the app every 5\,s and evaluates the HLD 12
alert rules (Table~\ref{tab:alerts}), and Grafana opens on a dashboard of the same numbers. The ops page shows the
store view: problems first (stalled orders, bags waiting too long for a rider, items out of stock), then eleven KPI
cards, each beside its target and marked met or not met.

\begin{table}[ht]
  \centering
  \footnotesize
  \begin{tabular}{@{}>{\raggedright\arraybackslash}p{0.27\linewidth}>{\raggedright\arraybackslash}p{0.68\linewidth}@{}}
    \toprule
    Alert & Fires when \\
    \midrule
    Failed holds spike & More than 20\% of checkout holds fail over 5 minutes \\
    Event lag & The oldest unsent event is older than 30\,s for a minute \\
    Slow reads & The p99 of any read route is above the 100\,ms target \\
    Store offline & A store has published nothing for 15 minutes \\
    Cache and database disagree & Waits for the stock audit (task 1.8); not wired up \\
    \bottomrule
  \end{tabular}
  \caption{HLD 12 alert rules, in \texttt{monitoring/alerts.yml}. The thresholds are first guesses to tune with real
    traffic.}
  \label{tab:alerts}
\end{table}

\paragraph{Load test.} \texttt{scripts/load\_test.py} sends the HLD's mix: about 90\% reads and 10\% checkout holds,
a fifth of them paid and the rest released. Eight clients ran for 60\,s against one laptop, with the rate limit off
(Table~\ref{tab:load}). Reads met the 100\,ms target with room to spare. No call failed: the 189 refused holds are
the store correctly running out of stock. The outbox lag stayed at 0.2\,s, one relay pass.

\begin{table}[ht]
  \centering
  \small
  \begin{tabular}{@{}lrrrr@{}}
    \toprule
    Endpoint & Calls & Per second & p50 (ms) & p99 (ms) \\
    \midrule
    Availability (read)   & 8{,}494 & 141.6 & 18.2 &  89.9 \\
    Serviceability (read) & 8{,}640 & 144.0 & 10.7 &  26.1 \\
    Hold                  & 1{,}918 &  32.0 & 22.0 & 142.9 \\
    Payment               &     313 &   5.2 & 36.5 & 121.7 \\
    Release               & 1{,}416 &  23.6 & 18.1 & 142.8 \\
    \bottomrule
  \end{tabular}
  \caption{Load test: 8 clients for 60\,s, SQLite on one laptop, main at b39692e. No 5xx; 189 holds refused for
    lack of stock.}
  \label{tab:load}
\end{table}

What the numbers say, and what they do not:
\begin{itemize}
  \item The cache served 51\% of availability reads. This mix sends one stock write for about every two availability
    reads, and each write clears the store's entry; Member 1's benchmark at the HLD's 13 reads per write measured
    92\%.
  \item Writes are slower than reads because SQLite lets one writer in at a time. Under an earlier load test the
    picker's \texttt{next-task}, which claims an order and so writes, reached a p99 of 250--500\,ms. PostgreSQL's row
    locks would let writes to different stores proceed side by side.
  \item This is one laptop, not the HLD's cluster: the numbers show where the time goes, not what production would
    do.
\end{itemize}

\section{Trade-off: a transactional outbox instead of Kafka (HLD 13)}
\label{sec:m4-tradeoff}

The HLD sends events through Kafka. The prototype writes each event to a table in the same transaction as the change
that caused it, and a relay inside the app delivers it.

\textbf{What it gains.} An event exists if and only if its change committed: no event announces a hold that was
rolled back, and no committed sale goes unannounced. That is the hard part of using Kafka well, and the outbox gets
it from one database transaction. It needs no broker to run or operate, and under load the lag stayed at one relay
pass (0.2\,s).

\textbf{What it costs.}
\begin{itemize}
  \item \textbf{At least once, not exactly once.} A sink that fails gets the batch again, so consumers must ignore an
    event id they have seen. The edge-node experiment checks both cases (\texttt{tests/test\_edge\_node.py}). With
    the link down, 7 events queued; once it was back, all 7 arrived in order, once each. A crash after sending
    delivered a batch twice, and the receiver kept each event once.
  \item \textbf{One relay, inside the app.} If the app is down, nothing is relayed; but nothing changes either, and
    the backlog drains on restart.
  \item \textbf{One database.} The outbox shares SQLite's single writer, and \texttt{event\_log} is never pruned.
    Replay for a new consumer means reading the table, not rewinding a topic.
\end{itemize}

\textbf{When to switch.} A Kafka producer is one more sink (\texttt{events.add\_sink}), so the outbox would stay as
the reliable hand-off and Kafka would carry events to other services. That is worth it once several services
consume events or the stores move to edge nodes.

\section{Limits and open items}

\begin{table}[!htb]
  \centering
  \footnotesize
  \begin{tabular}{@{}>{\raggedright\arraybackslash}p{0.27\linewidth}>{\raggedright\arraybackslash}p{0.36\linewidth}
                  >{\raggedright\arraybackslash}p{0.31\linewidth}@{}}
    \toprule
    Limit & What it means & Next step \\
    \midrule
    No separate gateway, no TLS & The app does the gateway's work itself & A gateway in front, TLS everywhere \\
    \addlinespace[3pt]
    One shared signing key (HS256) & Whoever can check a token can also make one & Asymmetric signing, needed
      anyway for edge nodes \\
    \addlinespace[3pt]
    Tokens cannot be revoked & A stolen token works until it expires (30 minutes) & Short tokens with refresh, or a
      deny list \\
    \addlinespace[3pt]
    Rate limit per process and per address & Several instances or a proxy share or split the buckets & Move it to
      the gateway or Redis \\
    \addlinespace[3pt]
    Audit log never pruned & It grows with every staff call & Keep a set period, then archive \\
    \addlinespace[3pt]
    Zones rebuilt by hand; one store per cell & A busy store's outer customers are not passed to a neighbour & Let a
      cell list a second store \\
    \addlinespace[3pt]
    Demand is generated & The map and the site suggestions run on scripted orders & Real order history \\
    \addlinespace[3pt]
    No Alertmanager & Alerts show only on Prometheus's page & Route them to the on-call channel \\
    \bottomrule
  \end{tabular}
  \caption{Limits of the platform pieces and what would close them.}
  \label{tab:m4-limits}
\end{table}
